\documentclass[11pt,a4paper]{article}

\usepackage{../shared/cathedral-preamble}

\title{%
  The Glass Bridge: Arithmetic Reduction of the Riemann\\
  Hypothesis via the Ramanujan--Rank-1 Decomposition%
}
\author{
  Jason Robert Gochanour\,\orcidlink{0000-0002-2862-526X}
}
\date{June 2026}

\begin{document}

\maketitle

\noindent
\textbf{MSC 2020:} 11M26 (primary), 11A25, 15A18, 68V15 (secondary).\\
\textbf{Keywords:} Riemann Hypothesis, Nyman--Beurling criterion,
Ramanujan matrix, Smith normal form, Sherman--Morrison, GCD lattice,
formal verification, Lean~4.

\begin{abstract}
We prove in Lean~4 (zero \lean{sorry}) that the Vasyunin Gram matrix
$G$ of the B\'aez-Duarte basis decomposes as
$G = R + \frac{1}{4}\mathbf{1}\mathbf{1}^\top$,
where $R_{jk} = \gcd(j,k)^2/(12jk)$ is the Ramanujan matrix and
$\mathbf{1}\mathbf{1}^\top$ is the all-ones rank-1 matrix. This
\emph{Glass Bridge} identity transforms the Nyman--Beurling distance
$d_N^2 = 1 - b^\top G^{-1} b$ into an explicit formula involving
only the inverse of~$R$. Via the Sherman--Morrison inversion
(machine-verified, zero axioms), the Riemann Hypothesis is equivalent
to the purely arithmetic condition $b^\top R^{-1} b \to 1$.
The Ramanujan matrix~$R$ admits a Smith normal form
diagonalization through the Jordan totient~$J_2$, reducing~RH
to the spectral behavior of a GCD lattice operator. We provide
numerical evidence through $N = 55{,}440$ using double-double
precision Gram matrices computed on GPU, and show that the
condition number explosion $\kappa(G) \sim N$ is entirely
an artifact of the rank-1 shift, while $\kappa(R) \sim \sqrt{N}$
and the dark crystal $\kappa(G^{(2)}) \sim 3$.
The Glass Bridge is also the mathematical foundation of the
\emph{anomaly separation} (Insight~E in the companion
physics paper~\cite{cathedral-physics}): the decomposition
$G = A + H$ into arithmetic and archimedean components that
underpins anomaly-gated precision engineering.
\end{abstract}

\tableofcontents

% ================================================================
\section{Introduction}
\label{sec:intro}
% ================================================================

\subsection{From Functional Analysis to Arithmetic}

The Nyman--Beurling criterion \cite{nyman1950,beurling1955}
reformulates the Riemann Hypothesis as a function-approximation
problem in~$L^2(0,1)$. In the B\'aez-Duarte basis
\cite{baezduarte2003}, define
\[
  d_N^2 = \inf_{v \in \RR^{N-1}} \int_0^1
    \biggl(1 - \sum_{k=1}^{N-1} v_k \fract{1/(kx)}\biggr)^2 dx.
\]
The Riemann Hypothesis holds if and only if $d_N^2 \to 0$.

The standard approach treats this as a continuous optimization
problem: compute the Gram matrix~$G$, invert it, and check
whether $b^\top G^{-1} b \to 1$. But the Gram matrix mixes
continuous geometry (fractional-part integrals) with discrete
arithmetic (GCD structure), obscuring the mechanism.

The Glass Bridge isolates the two. We prove that $G$ splits
exactly into a pure arithmetic component~$R$ and a trivial
rank-1 constant. The Riemann Hypothesis then reduces to a
statement about the inverse of a matrix whose entries are
$\gcd(j,k)^2/(12jk)$---pure number theory.

\subsection{Main Results}

\begin{theorem}[The Glass Bridge Identity, Machine-Verified]
\label{thm:glass}
For all $j, k \geq 1$,
\[
  G_{jk} = \frac{\gcd(j,k)^2}{12\,j\,k} + \frac{1}{4}.
\]
Equivalently, $G = R + \frac{1}{4}\mathbf{1}\mathbf{1}^\top$,
where $R_{jk} = \gcd(j,k)^2/(12jk)$ is the Ramanujan matrix.
\end{theorem}

\begin{theorem}[Arithmetic Reduction of RH, Machine-Verified]
\label{thm:arith-rh}
The Riemann Hypothesis is equivalent to
\[
  b^\top R^{-1} b - \frac{(b^\top R^{-1} \mathbf{1})^2}
    {4 + \mathbf{1}^\top R^{-1} \mathbf{1}} \;\longrightarrow\; 1
  \quad\text{as } N \to \infty,
\]
where $b_k = 1/2 - 1/(2k)$ and $R$ is the
$(N{-}1) \times (N{-}1)$ Ramanujan matrix.
\end{theorem}

\noindent
Both theorems compile with \textbf{zero sorry} and \textbf{zero
custom axioms}. No PNT input is needed---this is a pure algebraic
reformulation.

% ================================================================
\section{The Glass Bridge Identity}
\label{sec:identity}
% ================================================================

\subsection{The Ramanujan Inner Product}

\begin{theorem}[Machine-Verified, Zero Sorry]
\label{thm:ramanujan-ip}
For $j, k \geq 1$,
\[
  \int_0^1 B_1(\fract{jt})\,B_1(\fract{kt})\,dt
  = \frac{\gcd(j,k)^2}{12\,j\,k},
\]
where $B_1(x) = x - 1/2$ is the first Bernoulli polynomial.
\end{theorem}

\begin{proof}[Proof sketch]
Write $B_1(\fract{jt}) = \fract{jt} - 1/2$. The integral
$\int_0^1 \fract{jt}\fract{kt}\,dt$ computes as a
Ramanujan-type GCD sum via the Chinese Remainder Theorem
permutation on divisor residues. The full 170-line proof is in
\lean{Spectral/RamanujanInnerProduct.lean}.
\end{proof}

\subsection{From Sawtooth to Fractional Part}

The Gram matrix uses fractional parts
$\fract{jt} = B_1(\fract{jt}) + 1/2$. Expanding:
\begin{align*}
  G_{jk} &= \int_0^1 \fract{jt}\,\fract{kt}\,dt \\
  &= \int_0^1 \bigl(B_1(\fract{jt}) + \tfrac{1}{2}\bigr)
     \bigl(B_1(\fract{kt}) + \tfrac{1}{2}\bigr)\,dt \\
  &= \underbrace{\int_0^1 B_1(\fract{jt})B_1(\fract{kt})\,dt}
     _{= \gcd(j,k)^2/(12jk)}
   + \frac{1}{2}\underbrace{\int_0^1 B_1(\fract{jt})\,dt}_{= 0}
   + \frac{1}{2}\underbrace{\int_0^1 B_1(\fract{kt})\,dt}_{= 0}
   + \frac{1}{4}.
\end{align*}
The cross terms vanish because $\int_0^1 B_1(\fract{nt})\,dt = 0$
for all $n \geq 1$ (the sawtooth function has mean zero on each
period). This yields:
\[
  G_{jk} = \frac{\gcd(j,k)^2}{12jk} + \frac{1}{4}
  = R_{jk} + \frac{1}{4}.
\]

\begin{remark}
In matrix form, $G = R + \frac{1}{4}\mathbf{1}\mathbf{1}^\top$.
The constant $1/4$ arises from $\int_0^1 (1/2)^2\,dt = 1/4$---the
squared mean of the fractional-part function. This is the
\emph{only} source of the rank-1 perturbation.
\end{remark}

% ================================================================
\section{The Quadratic Form Decomposition}
\label{sec:quadform}
% ================================================================

\begin{theorem}[Machine-Verified, Zero Sorry]
\label{thm:quadform}
For any vector $v \in \RR^{N-1}$,
\[
  v^\top G v = v^\top R v
    + \frac{1}{4}\bigl(\textstyle\sum_k v_k\bigr)^2.
\]
\end{theorem}

\begin{proof}
By the Glass Bridge identity,
$\sum_{j,k} v_j G_{jk} v_k
= \sum_{j,k} v_j R_{jk} v_k + \frac{1}{4} \sum_{j,k} v_j v_k
= v^\top R v + \frac{1}{4}(\sum_k v_k)^2$.
\end{proof}

This splits the Nyman--Beurling distance into three independent
terms:
\begin{equation}
\label{eq:three-terms}
  d_N^2 = 1 - 2b^\top v + v^\top R v
    + \tfrac{1}{4}(\textstyle\sum_k v_k)^2.
\end{equation}

\begin{itemize}
  \item \textbf{The Ramanujan residual} $v^\top R v$: positive
    semi-definite, encodes GCD arithmetic. For the M\"obius witness
    $v_k = \mu(k)/k$, this converges to $1/(2\pi^2)$ via the Euler
    product (Section~\ref{sec:euler}).
  \item \textbf{The rank-1 noise}
    $\frac{1}{4}(\sum_k v_k)^2$: killed by PNT, since
    $\sum \mu(k)/k \to 0$.
  \item \textbf{The linear term} $-2b^\top v$: the BD integrals.
\end{itemize}

% ================================================================
\section{Sherman--Morrison Inversion}
\label{sec:sm}
% ================================================================

Since $G = R + \frac{1}{4}\mathbf{1}\mathbf{1}^\top$ is a rank-1
update of~$R$, the Sherman--Morrison formula gives:

\begin{theorem}[Machine-Verified, Zero Sorry, Zero Axioms]
\label{thm:sm}
If $R$ is invertible and $4 + \sigma \neq 0$ where
$\sigma = \mathbf{1}^\top R^{-1} \mathbf{1}$, then
\[
  G^{-1} = R^{-1}
    - \frac{R^{-1} \mathbf{1}\,\mathbf{1}^\top R^{-1}}
      {4 + \sigma}.
\]
\end{theorem}

\noindent
The Nyman--Beurling distance becomes:
\begin{align}
  d_N^2 &= 1 - b^\top G^{-1} b \notag \\
  &= 1 - b^\top R^{-1} b
    + \frac{(b^\top R^{-1} \mathbf{1})^2}{4 + \sigma}.
  \label{eq:glass-distance}
\end{align}

\begin{corollary}[The Glass Distance Formula]
The Riemann Hypothesis is equivalent to
\begin{equation}
\label{eq:rh-glass}
  b^\top R^{-1} b
    - \frac{(b^\top R^{-1} \mathbf{1})^2}{4 + \sigma}
  \;\longrightarrow\; 1.
\end{equation}
\end{corollary}

\noindent
This is a statement about the inverse of $R$---a matrix whose
entries are $\gcd(j,k)^2/(12jk)$. No fractional parts.
No integrals. No complex analysis. Pure GCD arithmetic.

% ================================================================
\section{The Smith Normal Form}
\label{sec:smith}
% ================================================================

\subsection{Smith Diagonalization of $R$}

The Ramanujan matrix $R$ admits a Smith-type factorization
through the Euler totient tower.

\begin{theorem}
\label{thm:smith}
Let $\Phi$ be the matrix with $\Phi_{nd} = [d \mid n]$, and
$D = \operatorname{diag}(1, 1/2, 1/3, \ldots)$. Then
\[
  R = \frac{1}{12} D^{-1} \Phi\,
    \operatorname{diag}(J_2(1), J_2(2), \ldots)\,
    \Phi^\top D^{-1},
\]
where $J_2(d) = d^2 \prod_{p \mid d}(1 - 1/p^2)$ is the
Jordan totient function.
\end{theorem}

\subsection{Implications for $R^{-1}$}

The Smith decomposition gives
$R^{-1} = 12\,D\,(\Phi^\top)^{-1}
\operatorname{diag}(1/J_2(d))\,\Phi^{-1} D$.
Since $(\Phi^{-1})_{dk} = \mu(d/k)\cdot[k \mid d]$
(M\"obius inversion), the entries of $R^{-1}$ are explicit
arithmetic functions. The spectral behavior of $R^{-1}$ is
controlled by the Jordan totient, which grows as
$J_2(d) \sim (6/\pi^2) d^2$.

% ================================================================
\section{The Euler Product: $v^\top R v \to 1/(2\pi^2)$}
\label{sec:euler}
% ================================================================

For the M\"obius witness $v_k = \mu(k)/k$, the Smith decomposition
gives
\[
  v^\top R v = \frac{1}{12} \sum_d J_2(d)\,y_d^2,
\]
where $y_d = \sum_{k \leq N,\, d \mid k} \mu(k)/k^2$.
Writing $k = dm$ with $\gcd(d,m) = 1$:
\[
  y_d = \frac{\mu(d)}{d^2} \cdot \frac{6}{\pi^2}
  \prod_{p \mid d} \frac{p^2}{p^2 - 1}
  \quad \text{as } N \to \infty.
\]

The Euler product evaluates to:
\begin{align*}
  v^\top R v &\;\longrightarrow\;
  \frac{3}{\pi^4} \prod_p \biggl[1 + \frac{1}{p^2 - 1}\biggr]
  = \frac{3}{\pi^4} \prod_p \frac{p^2}{p^2 - 1}
  = \frac{3}{\pi^4} \cdot \zeta(2)
  = \frac{3}{\pi^4} \cdot \frac{\pi^2}{6} \\
  &= \frac{1}{2\pi^2} \approx 0.050661.
\end{align*}

\noindent
Numerically verified to $2.76 \times 10^{-6}$ at $N = 10{,}000$.

\begin{remark}
Since $v^\top R v \to 1/(2\pi^2) > 0$, the simple M\"obius
witness $v_k = \mu(k)/k$ does \emph{not} make $d_N^2 \to 0$.
The optimal witness $v^* = G^{-1} b$ achieves $d_N^2 = 1 - b^\top G^{-1} b$,
and RH asks precisely whether this approaches zero.
Through the Glass Bridge, that question lives entirely in $R^{-1}$.
\end{remark}

% ================================================================
\section{The Von Mangoldt Bridge}
\label{sec:vonmangoldt}
% ================================================================

\subsection{The Bridge Identity}

The mean vector in the Smith basis reveals a striking connection
to the von Mangoldt function.

\begin{theorem}[Machine-Verified, Zero Sorry]
\label{thm:vm-bridge}
Define $c_d = \sum_{k \mid d} \mu(d/k)(\log k + 1 - \gamma)$.
Then
\[
  c_d = \Lambda(d) + (1 - \gamma)\cdot[d = 1],
\]
where $\Lambda$ is the von Mangoldt function and $\gamma$ is the
Euler--Mascheroni constant.
\end{theorem}

\noindent
This identity, proved in \lean{VonMangoldtBridge.lean}, shows that
the Gram matrix's mean vector, when rotated into the Smith
eigenbasis, projects entirely onto the prime powers. The arithmetic
of the Riemann Hypothesis is concentrated on $\Lambda$.

% ================================================================
\section{Condition Number Hierarchy}
\label{sec:condition}
% ================================================================

The Glass Bridge explains the condition number hierarchy observed
in numerical experiments:

\begin{center}
\begin{tabular}{@{}rccc@{}}
\toprule
$N$ & $\kappa(G)$ & $\kappa(R)$ & $\kappa(G^{(2)})$ \\
\midrule
10 & 109 & 8.6 & 2.5 \\
20 & 260 & 12.6 & 2.8 \\
50 & 751 & 18.0 & 3.2 \\
80 & 1{,}298 & 21.2 & 3.4 \\
\bottomrule
\end{tabular}
\end{center}

\noindent
The full Gram matrix $G$ has $\kappa \sim N$ (ill-conditioned).
After removing the rank-1 shift, $R$ has $\kappa \sim \sqrt{N}$
(mild growth). The dark Gram matrix $G^{(2)}_{jk} =
\gcd(j,k)^4/(180\,j^2 k^2)$ remains at $\kappa \sim 3$
(nearly perfectly conditioned). The chaos of~$G$ is entirely
an artifact of adding $\frac{1}{4}$ to every entry.

% ================================================================
\section{Numerical Evidence}
\label{sec:numerics}
% ================================================================

Computation of $v^\top G v$ for the M\"obius logCutoff witness
$v_k = \mu(k)/\log N$, using DD-precision HPDF Gram matrices
on GPU:

\begin{center}
\begin{tabular}{@{}rccc@{}}
\toprule
$N$ & $v^\top G v$ & $1 - v^\top G v$ & $d_N^2$ \\
\midrule
120 & 0.5577 & 0.4423 & 0.215 \\
720 & 0.8692 & 0.1308 & 0.348 \\
2{,}520 & 0.9389 & 0.0611 & 0.342 \\
5{,}040 & 0.8516 & 0.1484 & 0.222 \\
\bottomrule
\end{tabular}
\end{center}

\noindent
All values satisfy $v^\top G v < 1$, consistent with the
overcancellation hypothesis established in the companion
paper~\cite{overcancellation}.

% ================================================================
\section{Discussion}
\label{sec:discussion}
% ================================================================

\subsection{The Arithmetic Core of RH}

The Glass Bridge reveals that the difficulty of the Riemann
Hypothesis is not in the continuous geometry of $L^2(0,1)$,
nor in the complex analysis of $\zeta(s)$. It is in the
arithmetic structure of the GCD lattice: the matrix
$R_{jk} = \gcd(j,k)^2/(12jk)$ and its inverse.

Through the Smith Normal Form, $R^{-1}$ is controlled by
the Jordan totient $J_2(d) = d^2 \prod_{p|d}(1-1/p^2)$,
and the mean vector projects onto the von Mangoldt function.
Every component of the Nyman--Beurling distance formula is
a classical arithmetic function.

\subsection{Open Questions}

\begin{enumerate}
  \item \textbf{Spectral asymptotics of $R^{-1}$.} Determine
    the growth rate of $b^\top R^{-1} b$ as $N \to \infty$
    to verify~\eqref{eq:rh-glass}.
  \item \textbf{GOE statistics.} Numerical evidence suggests
    the Gram matrix eigenvalues follow Gaussian Orthogonal
    Ensemble statistics. Does the Smith basis explain this?
  \item \textbf{Dark crystal universality.} The dark Gram matrix
    $G^{(2)}$ has condition number $\kappa \sim 3$ for all
    tested~$N$. Is this universal?
  \item \textbf{Connection to $L$-functions.} Does the Glass
    Bridge generalize to Dirichlet $L$-functions, replacing
    $\mu$ with twisted M\"obius functions?
  \item \textbf{The Glass Bridge as anomaly separator.} The
    decomposition $G = R + \frac{1}{4}\mathbf{1}\mathbf{1}^\top$
    is an exact instance of the arithmetic/archimedean anomaly
    separation. Does the Torus Projection
    $T^\infty = \prod_p S^1_p$ further decompose $R$ into
    per-prime strata with independent spectral behavior?
\end{enumerate}

% ================================================================
\section{Formalization Architecture}
\label{sec:lean}
% ================================================================

\begin{center}
\small
\begin{tabular}{@{}llcc@{}}
\toprule
\textbf{Component} & \textbf{File} & \textbf{Sorry} & \textbf{Axioms} \\
\midrule
Ramanujan inner product & \lean{RamanujanInnerProduct.lean} & 0 & 0 \\
Glass Bridge identity & \lean{RamanujanBridge.lean} & 0 & 0 \\
Quadratic form split & \lean{RamanujanBridge.lean} & 0 & 0 \\
PSD of $R$ & \lean{RamanujanBridge.lean} & 0 & 0 \\
Sherman--Morrison & \lean{ShermanMorrison.lean} & 0 & 0 \\
Von Mangoldt bridge & \lean{VonMangoldtBridge.lean} & 0 & 0 \\
Smith witness & \lean{SmithWitness.lean} & 0 & 0 \\
\textbf{Glass distance} & \lean{GramBridge.lean} & \textbf{0} & \textbf{0} \\
\bottomrule
\end{tabular}
\end{center}

\noindent
The entire Glass Bridge pathway compiles with \textbf{zero sorry}
and \textbf{zero custom axioms}. This is a pure algebraic
reduction requiring no PNT input.

% ================================================================
\section*{Acknowledgments}
% ================================================================

The author thanks Claude (Anthropic) and Gemini (Google DeepMind)
for computational and theoretical contributions to the Glass Bridge
discovery and formalization campaign (March--June 2026).

% ================================================================
\begin{thebibliography}{99}
% ================================================================

\bibitem{nyman1950}
B.~Nyman,
``On the one-dimensional translation group and semi-group in certain
function spaces,''
Thesis, Uppsala, 1950.

\bibitem{beurling1955}
A.~Beurling,
``A closure problem related to the Riemann zeta function,''
\emph{Proc.\ Nat.\ Acad.\ Sci.\ USA} \textbf{41} (1955), 312--314.

\bibitem{baezduarte2003}
L.~B\'aez-Duarte,
``A strengthening of the Nyman--Beurling criterion for the Riemann
hypothesis,''
\emph{Atti Accad.\ Naz.\ Lincei} \textbf{14} (2003), 5--11.

\bibitem{vasyunin1995}
V.~I.~Vasyunin,
``On a biorthogonal system associated with the Riemann hypothesis,''
\emph{Algebra i Analiz} \textbf{7} (1995), 118--135.

\bibitem{ramanujan1918}
S.~Ramanujan,
``On certain trigonometrical sums and their applications in the
theory of numbers,''
\emph{Trans.\ Cambridge Philos.\ Soc.} \textbf{22} (1918), 259--276.

\bibitem{overcancellation}
J.~R.~Gochanour,
``The overcancellation path to the Riemann Hypothesis:
a formalized conditional proof in Lean~4,''
2026.

\bibitem{cathedral-main}
J.~R.~Gochanour,
``A formal reduction of the Riemann Hypothesis to two axioms
via the Nyman--Beurling criterion in Lean~4,''
2026.

\bibitem{cathedral-physics}
J.~R.~Gochanour,
``The Physics of the Primes: A Dictionary Between the Cathedral
Proof and Quantum Mechanics,''
2026.

\bibitem{mathlib}
The Mathlib Community,
``Mathlib: the mathematics library for Lean~4,''
\url{https://github.com/leanprover-community/mathlib4}.

\end{thebibliography}

\vfill
\begin{center}
\itshape
The difficulty of the Riemann Hypothesis is not in the chaos.\\
It is in the arithmetic. And the arithmetic is
$\gcd(j,k)^2/(12jk)$.
\end{center}

\end{document}
